%!TEX root = ../thesis.tex
% ******************************* Thesis Appendix B ********************************

\chapter{Additional Results for Sampling-Based Bayesian Inference}

\label{app:gp_bnn}

\section{Additional Proofs}

% \subsection{Proof of \Cref{prop:spectral_basis}}
% \label{app:courant_in_input_space}

% Since the preceding seminorm is induced by a projection within the RKHS, the next step is to understand what kind of functions the corresponding subspace contains.
% To get a qualitative view, the basic idea is to lift the Courant--Fischer characterization of eigenvalues and eigenvectors to the~RKHS.

% We will need the following variant of the Min-Max Theorem, which is useful for describing multiple eigenvectors at once---as solutions to a sequence of Raleigh-quotient-type optimization problems, where the next optimization problem is performed on a subspace which ensures its solution is orthogonal to all previous solutions.
% Mirroring the rest of this work, we consider eigenvalues in decreasing order, namely $\lambda_1 \geq .. \geq \lambda_N \geq 0$.

% \begin{result}
% \label{res:courant-fischer}
% Let $A$ be a positive semi-definite matrix.
% Then 
% \[
% \lambda_i &= \max_{\substack{u \in \R^N\takeaway\{{0}\}\\u^T u_j = 0, \forall j < i}} \frac{u^TAu}{u^Tu}
% \]
% where the eigenvectors $u_i$ are the respective maximizers.
% \end{result}

% \begin{proof}
% \textcite[Theorem 4.2.2]{horn12}, modified appropriately to handle eigenvalues in decreasing order, and where we choose $i_1,..,i_k = 1,..,N-i+1$.
% \end{proof}

% We will prove the main claim below in two stages: first, on representer weight space $R_{k,{x}}$, then on the full RKHS $H_k$.
% We begin with the former.

% \begin{selflemma}
% \label{lem:courant-fischer-rkhs}
% We have 
% \[
% u^{(i)}(\.) &= \argmax_{u \in R_{k,{x}}} \cbr{\sum_{i=1}^N u(x_i)^2 : \norm{u}_{H_k} = 1, \innerprod[0]{u}{u^{(j)}}_{H_k} = 0, \forall j < i}.
% \]
% \end{selflemma}

% \begin{proof}
% Define ${w}_j = {U} \m\Lambda^{-1/2} u_j$, where we recall that $u_j$ are the respective eigenvectors of the kernel matrix.
% From the \Cref{res:courant-fischer} form of the Courant--Fischer Min-Max Theorem, the reproducing property, and the explicit form of the RKHS inner product on $R_{k,{x}}$ in the basis defined by the canonical basis functions, we can conclude that
% \allowdisplaybreaks
% \[
% \lambda_i &= \max_{\substack{u\in\R^N\takeaway\{{0}\}\\u^Tu_j = 0, \forall j<i}} \frac{u^T {K}_{{x}{x}} u}{u^Tu} = \max_{\substack{{w}\in\R^N\takeaway\{{0}\}\\{w}^T{K}_{{x}{x}}{w}_j = 0,\forall j<i}} \frac{{w}^T {K}_{{x}{x}}^2 {w}}{{w}^T {K}_{{x}{x}} {w}} 
% \\
% &= \max_{\substack{{w}\in\R^N\takeaway\{{0}\}\\{w}^T{K}_{{x}{x}}{w}_j = 0,\forall j<i}} \frac{\norm{{K}_{{x}{x}} {w}}_2^2}{{w}^T {K}_{{x}{x}} {w}} = \max_{\substack{{w}\in\R^N\takeaway\{{0}\}\\{w}^T{K}_{{x}{x}}{w}_j = 0,\forall j<i}} \frac{\norm{\sum_{i=1}^N w_i k(x_i,{x})}_2^2}{\norm{\sum_{i=1}^N w_i k(x_i,\.)}_{H_k}^2} 
% \\
% &= \max_{\substack{u \in R_{k,{x}}\takeaway\{0\}\\\innerprod[0]{u}{u^{(j)}}_{H_k} = 0, \forall j < i}} \frac{\norm{u({x})}_2^2}{\norm{u}_{H_k}^2} = \max_{\substack{u \in R_{k,{x}}\\\norm{u}_{H_k}=1\\\innerprod[0]{u}{u^{(j)}}_{H_k} = 0, \forall j < i}} \sum_{i=1}^N u(x_i)^2
% \]
% which gives the claim.
% \end{proof}

% Next, we use a Representer-Theorem-type orthogonality argument to extend the optimization problem to all of $H_k$.
% The argument will rely heavily on the Projection Theorem for Hilbert spaces: for an overview of the setting, see \textcite[Chapter 5]{lang12}.

% \begin{selfproposition}
% We have 
% \[
% u^{(i)}(\.) &= \argmax_{u \in H_k} \cbr{\sum_{i=1}^N u(x_i)^2 : \norm{u}_{H_k} = 1, \innerprod[0]{u}{u^{(j)}}_{H_k} = 0, \forall j < i}.
% \]
% \end{selfproposition}

% \begin{proof}
% Let $H_k^{(i)}$ be the orthogonal complement of the finite-dimensional subspace $\Span\{u^{(j)}, j=1,..,i-1\}$ in $H_k$ and note by feasibility that $u^{(i)} \in H_k^{(i)}$.
% Consider the decomposition of $u^{(i)}$ into its projection onto an orthogonal complement with respect to the finite-dimensional subspace $H_k^{(i)} \^ R_{k,{x}}$, namely
% \[
% u^{(i)}(\.) = u^\parallel(\.) + u^\orth(\.)
% \]
% where this notation is defined as $u^\parallel = \proj_{H_k^{(i)}\^R_{k,{x}}} u^{(i)}$ and $u^\orth = u^{(i)} - u^\parallel$.
% Observe that
% \[
% u^\orth(x_i) = \innerprod[0]{u^\orth}{k(x_i,\.)}_{H_k} = \innerprod[0]{\proj_{H_k^{(i)}} u^\orth}{k(x_i,\.)}_{H_k} = \innerprod[0]{u^\orth}{\proj_{H_k^{(i)}} k(x_i,\.)}_{H_k} = 0
% \]
% where the first equality follows from the reproducing property, the second equality follows because $u^\orth \in H_k^{(i)}$, the third equality follows because the projection is orthogonal and therefore self-adjoint, and last equality follows since $u^\orth$ by definition lives in the orthogonal complement of $H_k^{(i)}\^R_{k,{x}}$ within $H_k^{(i)}$, with the subspace inner product.
% Moreover, by the Projection Theorem and the fact that $H_k^{(i)}$ inherits its norm from $H_k$, note that 
% \[
% \norm[0]{u^{(i)}}_{H_k}^2 = \norm[0]{u^\parallel}_{H_k}^2 + \norm[0]{u^\orth}_{H_k}^2
% .
% \]
% We now argue by contradiction: suppose that $u^\orth \neq 0$.
% This means that $\norm[0]{u^\orth}_{H_k} > 0$.
% From the above, we have 
% \[
% \sum_{i=1}^N u^{(i)}(x_i)^2 &= \sum_{i=1}^N u^\parallel(x_i)^2
% &
% \norm[0]{u^\parallel}_{H_k} &< \norm[0]{u^{(i)}}_{H_k}
% .
% \]
% Define the function 
% \[
% u'(\.) = \frac{\norm[0]{u^{(i)}}_{H_k}}{\norm[0]{u^\parallel}_{H_k}} u^\parallel(\.)
% \]
% and note that $\norm{u'} = 1$, and $\innerprod[0]{u'}{u^{(j)}} = 0$ for all $j < i$ since $u^\parallel \in H_k^{(i)}$, which makes $u'$ a feasible solution to the optimization problem considered.
% At the same time, it satisfies 
% \[
% \sum_{i=1}^N u'(x_i)^2 = \frac{\norm[0]{u^{(i)}}_{H_k}^2}{\norm[0]{u^\parallel}_{H_k}^2}\sum_{i=1}^N u^\parallel(x_i)^2 > \sum_{i=1}^N u^\parallel(x_i)^2 = \sum_{i=1}^N u^{(i)}(x_i)^2
% \]
% which shows that $u^{(i)}$, contrary to its definition, is not actually optimal---contradiction.
% We conclude that $u^\orth = 0$, which means that $u^{(i)} \in H_k^{(i)} \^ R_{k,{x}}$, and the claim follows from \Cref{lem:courant-fischer-rkhs}.
% \end{proof}

% This means that the top eigenvalues can be understood as the largest possible squared sums  of canonical basis functions evaluated at the data, subject to RKHS-norm constraints.
% In situations where $k(x,\.)$ is continuous, non-negative, and decays as one moves away from the data, it is intuitively clear that such sums will be maximized by placing weight on canonical basis functions which cover as much of the data as possible.
% This mirrors the RKHS view of kernel principal component analysis.

\subsection{Proof of \Cref{remark:strong_duality}}
\label{app:proof_strong_duality}

\begin{proof}
That $\alpha^{\star}(b)$ minimises both $L$ and $L^{\star}$ can be established from the first order optimality conditions. Now, for the duality, observe that we can write $\min _{\alpha \in \mathbb{R}^n} L(\alpha)$ equivalently as the constrained optimisation problem
$$
\min _{u \in \mathbb{R}^n} \min _{\alpha \in \mathbb{R}^n} \frac{1}{2}\|u\|^2+\frac{\lambda}{2}\|\alpha\|_K^2 \quad \text { subject to } \quad u=K \alpha-b
$$


Observe that this is quadratic in both $u$ and $\alpha$. Introducing Lagrange multipliers $\beta \in \mathbb{R}^n$, in the form $\lambda \beta$, where we recall that $\lambda>0$, the solution of the above is equal to that of

$$
\min _{u \in \mathbb{R}^n} \min _{\alpha \in \mathbb{R}^n} \sup _{\beta \in \mathbb{R}^n} \frac{1}{2}\|u\|^2+\frac{\lambda}{2}\|\alpha\|_K^2+\lambda \beta^{\top}(b-K \alpha-u)
$$

This is a finite-dimensional quadratic problem, and thus we have strong duality (see, e.g., Examples 5.2.4 in Boyd and Vandenberghe, 2004). We can therefore exchange the order of the minimum operators and the supremum, yielding the again equivalent problem
$$
\sup _{\beta \in \mathbb{R}^n}\left\{\min _{u \in \mathbb{R}^n} \frac{1}{2}\|u\|^2-\lambda \beta^{\top} u\right\}+\left\{\min _{\alpha \in \mathbb{R}^n} \frac{\lambda}{2}\|\alpha\|_K^2-\lambda \beta^{\top} K \alpha\right\}+\lambda \beta^{\top} b
$$

Noting that the two inner minimisation problems are quadratic, we solve these analytically using the first order optimality conditions, that is $\alpha=\beta$ and $u=\lambda \beta$, to obtain that the above is equivalent to
$$
\sup _{\beta \in \mathbb{R}^n}-\lambda\left(\frac{1}{2}\|\beta\|_{K+\lambda I}^2-\beta^{\top} b\right)=-\lambda \min _{\beta \in \mathbb{R}^n} L^*(\beta) .
$$

The result follows by chaining the above equalities.
\end{proof}

\subsection{Proof of \Cref{prop:uniform_approximation}}
\label{app:proof_uniform_approximation}

To show the uniform approximation bound, we first need to define reproducing kernel Hilbert spaces (RKHSes). Let $\mathcal{H}$ be a Hilbert space of functions $\mathcal{X} \rightarrow \mathbb{R}$ with inner product $\langle\cdot, \cdot\rangle$ and corresponding norm $\|\cdot\|_{\mathcal{H}}$. We say $\mathcal{H}$ is the RKHS associated with a bounded kernel $k$ if the reproducing property holds:
\[
\forall x \in \mathcal{X}, \forall h \in \mathcal{H}, \quad\langle k(x, \cdot), h\rangle=h(x)
\]
That is, $k(x, \cdot)$ is the evaluation functional (at $x \in \mathcal{X}$ ) on $\mathcal{H}$. For observations $X$, let $\Phi: \mathcal{H} \rightarrow \mathbb{R}^n$ be the linear operator mapping $h \mapsto h(X)$, where $h(X)=\left(h\left(x_1\right), \ldots, h\left(x_n\right)\right)$. We will write $\Phi^*$ for the adjoint of $\Phi$, and observe that $K$ is the matrix of the operator $\Phi \Phi^*$ with respect to the standard basis.
\begin{proof}
First, observe that,
\[
\begin{aligned}
\|h_\alpha-h_{\alpha^{\prime}}\|_{\infty} & =\sup _{x \in \mathcal{X}}\left|h_\alpha(x)-h_{\alpha^{\prime}}(x)\right| \\
& =\sup _{x \in \mathcal{X}}\left|\left\langle k(x, \cdot), h_\alpha-h_{\alpha^{\prime}}\right\rangle\right| \\
& \leq \sup _{x \in \mathcal{X}}\|k(x, \cdot)\|_{\mathcal{H}}\|h_\alpha-h_{\alpha^{\prime}}\|_{\mathcal{H}} \\
& \leq \sqrt{\kappa}\|h_\alpha-h_{\alpha^{\prime}}\|_{\mathcal{H}}
\end{aligned}
\]
Now, observe that $h_\alpha=\Phi^* \alpha$ and $h_{\alpha^{\prime}}=\Phi^* \alpha^{\prime}$, and so we have the equalities
\[
\begin{aligned}
\|h_\alpha-h_{\alpha^{\prime}}\|_{\mathcal{H}}^2 & =\left\langle\Phi^*\left(\alpha-\alpha^{\prime}\right), \Phi^*\left(\alpha-\alpha^{\prime}\right)\right\rangle \\
& =\left\langle\alpha-\alpha^{\prime}, \Phi \Phi^*\left(\alpha-\alpha^{\prime}\right)\right\rangle \\
& =\|\alpha-\alpha^{\prime}\|_K^2
\end{aligned}
\]
where for the final equality, we note that $K$ is the matrix of the operator $\Phi \Phi^*$ with respect to the standard basis. Combining the above two displays yields the claim.
\end{proof}

\subsection{Proof of \Cref{prop:pathwise_dist}}
\label{app:proof_pathwise_dist}
We consider the generic linear system $H_\theta u=b$ and measure the RKHS distance between the initialisation $u_{\text {init }}$ and the solution $u=H_{\varphi}^{-1} b$ as $\|u_{\text {init }}-u\|_{H_{\varphi}}^2$. With $u_{\text {init }}=0$, which is standard \cite{gardner18,WangPGT2019exactgp,wu2024largescale}:
\[
\|u_{\text {init }}-u\|_{H_{\varphi}}^2=\|u\|_{H_{\varphi}}^2=u^{\top} H_{\varphi} u=b^{\top} H_{\varphi}^{-1} H_{\varphi} H_{\varphi}^{-1} b=b^{\top} H_{\varphi}^{-1} b \label{eq:RKHS_distance}
\]
Since $b$ is a random vector ($z$ in \Cref{eq:linear_system_trace_estimation} and $\xi$ in \Cref{eq:pathwise_trace_linear_system}), we analyse the expected squared distance
\[
\mathbb{E}\left[b^{\top} H_{\varphi}^{-1} b\right]=\mathbb{E}\left[\operatorname{tr}\left(b^{\top} H_{\varphi}^{-1} b\right)\right]=\mathbb{E}\left[\operatorname{tr}\left(b b^{\top} H_{\varphi}^{-1}\right)\right]=\operatorname{tr}\left(\mathbb{E}\left[b b^{\top}\right] H_{\varphi}^{-1}\right)
\]
For the standard estimator (6), we substitute $b:=z$ with $\mathbb{E}\left[z z^{\top}\right]=I$, yielding
\[
\mathbb{E}\left[\|u_{\text {init }}-u\|_{H_{\varphi}}^2\right]=\operatorname{tr}\left(\mathbb{E}\left[z z^{\top}\right] H_{\varphi}^{-1}\right)=\operatorname{tr}\left(I H_{\varphi}^{-1}\right)=\operatorname{tr}\left(H_{\varphi}^{-1}\right) . \label{eq:standard_dist_app}
\]
For the pathwise estimator (9), we substitute $b:=\xi$ with $\mathbb{E}\left[\xi \xi^{\top}\right]=H_{\varphi}$, yielding:
\[
\mathbb{E}\left[\|u_{\text {init }}-u\|_{H_{\varphi}}^2\right]=\operatorname{tr}\left(\mathbb{E}\left[\xi \xi^{\top}\right] H_{\varphi}^{-1}\right)=\operatorname{tr}\left(H_{\varphi} H_{\varphi}^{-1}\right)=\operatorname{tr}(I)=n . \label{eq:pathwise_dist_app}
\]

\subsection{Proof of \Cref{def:mackay_first_order_solution}}
\label{app:proof_mackay_first_order_solution}

\begin{proof}
Consider the first-order derivative of \Cref{eq:vi_fixed_cov_A} in \Cref{def:fixed_cov_elbo} w.r.t $A$:
\[
\frac{\partial \mathcal{M}(\theta_q; A)}{\partial A} = -\frac{1}{2} \left[ \nabla \| \theta_q \|_{A}^2 + \nabla_A \log \det (A + \Phi \Sigma^{-1} \Phi^T)  - \nabla_A \log \det A\right] \label{eq:mackay_first_order_derivative}
\]
Taking derivatives and setting to zero, we obtain the following condition:
\[
    \theta_q \theta_q^T = \left(I-\left(I+A^{-1} (\Phi \Sigma^{-1} \Phi^T)\right)^{-1}\right) A^{-1}
\]
Post-multiplying with $A$ and applying the push-forward identity, we obtain:
\[
    \theta_q \theta_q^T A=(\Phi \Sigma^{-1} \Phi^T)(A+\Phi \Sigma^{-1} \Phi^T)^{-1}
\]
For the above to hold, the traces on both sides must be equal. Therefore we have:
\[
    \|\theta_q\|_A^2&=\Tr\left\{\theta_q \theta_q^T A\right\}=\Tr\left\{(\Phi \Sigma^{-1} \Phi^T)(A+\Phi \Sigma^{-1} \Phi^T)^{-1}\right\} \\
    &= \Tr\left\{(\Phi \Sigma^{-1} \Phi^T)A_{\star}^{-1}\right\}
\]
which is the stated first order optimality condition, up to a cyclic permutation.
\end{proof}

\subsection{Proof of \Cref{def:matheron_rule_linear_model}}
\label{app:proof_matheron_rule_linear_model}
\begin{proof}
    We begin from \Cref{eq:pathwise_conditioning_gp} and plug in the form of the kernel $K_{XX} = \Phi A^{-1} \Phi^T$ and $K_{(\cdot)X} = \phi(\cdot) A^{-1} \Phi^T$. We also have $f(X) = \Phi \theta$ from \Cref{eq:linear-model-functional-space} and $(f \mid Y)(\cdot) = \phi(\cdot) \theta$ from \Cref{eq:linear-model-posterior-functional-space}. We then have:
    \[
        \phi(\cdot) (\theta \mid \rmY) &= \phi(\cdot) \theta + \phi(\cdot) A^{-1} \Phi^T\left(\Phi A^{-1} \Phi^T + \Sigma\right)^{-1}\left(Y-\Phi \theta - \varepsilon\right) \nonumber
    \]
    Removing the multiplicative factor of $\phi(\cdot)$ from both sides, we obtain \Cref{eq:matheron_rule_linear_model_1}.
    \[
        (\theta \mid \rmY) = \theta_0 + A^{-1} \Phi^T (\Phi A^{-1} \Phi^T + \Sigma)^{-1} (\rmY - \Phi \theta_0 - \varepsilon). \nonumber
    \]
    Instead, beginning from \Cref{eq:pathwise_conditioning_gp_2}, we can write the zero-mean parameter sample equivalently.
\end{proof}

\subsection{Proof of Definition \ref{def:zero_mean_posterior_samples}}
\label{app:proof_zero_mean_posterior_samples}

\begin{proof}
    We manipulate \Cref{eq:matheron_rule_linear_model_2} to obtain:
    \begin{align*}
        \bar{\theta}_{\star} &= \theta_0 + A^{-1} \Phi^T\left(\Phi A^{-1} \Phi^T+\Sigma\right)^{-1}\left(\Phi \theta_0 - \varepsilon\right)\\
        &= \theta_0 + A^{-1} \Phi^T \Sigma^{-1} (\Phi A^{-1} \Phi^T \Sigma^{-1} )^{-1} \left(\Phi \theta_0 - \varepsilon\right) \tag{Taking $B$ outside} \\
        &= \theta_0 + A^{-1} (I + \Phi ^T \Sigma^{-1} \Phi A^{-1})^{-1} \Phi^T \Sigma^{-1} \left(\Phi \theta_0 - \varepsilon\right) \tag{Morrison-Sherman formula} \\
        &= \theta_0 + (A + \Phi ^T \Sigma^{-1} \Phi)^{-1} \Phi^T \Sigma^{-1} \left(\Phi \theta_0 - \varepsilon\right) \tag{Taking $A^{-1}$ inside}\\
        &= A_{\star}^{-1} A_{\star} \theta_0 + A_{\star}^{-1} \Phi^T \Sigma^{-1} \left(\Phi \theta_0 - \varepsilon\right) \\
        &= A_{\star}^{-1} ( (A_{\star} - \Phi \Sigma^{-1} \Phi) \theta_0 + \Phi \Sigma^{-1} \varepsilon) \\
        &= A_{\star}^{-1} (A \theta_0 + \Phi \Sigma^{-1} \varepsilon), \quad \theta_0 \sim \mathcal{N}(0, A^{-1}), \varepsilon \sim \mathcal{N}(0, \Sigma).
    \end{align*}
    Now, comparing the form of \Cref{eq:linear-model-posterior-mean}, we can see that the term above is the posterior mean of a linear model with a prior specified as $\theta \sim \mathcal{N}(A \theta_0, A^{-1})$ and targets given by $\varepsilon$ instead of $\rmY$. Therefore, we can write down a maximum a-posteriori loss function similar to \Cref{eq:linear-model-posterior-log-density} for this linear model instead as:
    \[
    L(\theta) = \frac{1}{2} \| \Phi \theta - \varepsilon\|^2_{\Sigma^{-1}} + \frac{1}{2} \| \theta - \theta_0 \|^2_{A}
    \]
\end{proof}

\subsection{Proof of Definition \ref{def:low_variance_objective_zero_mean_posterior_samples}}
\label{app:proof_low_variance_objective_zero_mean_posterior_samples}
\begin{proof}

    The losses $L$ and $L'$ are strictly convex, thus to confirm they have the same unique minimum, it suffices to consider the respective first order optimality conditions, $\nabla L(\bar{\theta}) = 0$ and $\nabla L'(\bar{\theta}') = 0$. We have,
\begin{equation}
  \nabla L (\bar{\theta}) = \Phi^T\Sigma^{-1} (\Phi \bar{\theta} - \varepsilon) + A(\bar{\theta} - \theta^0),
\end{equation}
and 
\begin{align}
  \nabla L' (\bar{\theta}') &= \Phi^T\Sigma^{-1}\Phi \bar{\theta}' + A(\bar{\theta}' - A^{-1}\Phi^T\Sigma^{-1}\varepsilon - \theta^0) \\&= \Phi^T\Sigma^{-1} (\Phi \bar{\theta}' - \varepsilon) + A(\bar{\theta}' - \theta^0)
\end{align}
Thus $\bar{\theta} = \bar{\theta}'$ almost surely. Moreover, $L'(\theta)=L(\theta)+C$ for all $\theta$, for $C$ a constant independent of $\theta$.

To determine the distribution of $\bar{\theta}$, note that it is a linear transformation of zero-mean Gaussian random variables, and thus itself a zero-mean Gaussian random variable. Rearranging the first order optimality condition, we find that
\begin{equation}
  \bar{\theta} = A_{\star}^{-1}(\Phi^T\Sigma^{-1}\varepsilon + A\theta^0).
\end{equation}
Thus 
\begin{align}
  \E[\bar{\theta}\bar{\theta}^T] &= A_{\star}^{-1} \E[(\Phi^T\Sigma^{-1}\varepsilon + A\theta^0)(\Phi^T\Sigma^{-1}\varepsilon + A\theta^0)^T] A_{\star}^{-1} \\
  &= A_{\star}^{-1} \left(\Phi^T \Sigma^{-1} \E[\varepsilon\varepsilon^T] \Sigma^{-1} \Phi + A\E[\theta^0\theta^0] A + 2\Phi^T\Sigma^{-1}\E[\varepsilon(\theta^0)^T]A\right)A_{\star}^{-1} \\
  &= A_{\star}^{-1} (\Phi^T \Sigma^{-1} \Phi + A) A_{\star}^{-1} = A_{\star}^{-1} A_{\star} A_{\star}^{-1} = A_{\star}^{-1}.
\end{align}
And so $\bar{\theta} \sim \mathcal{N}(0, A_{\star}^{-1})$ as claimed.

\end{proof}

\subsection{Proof of \Cref{lem:gradient_variance_for_objectives}}
\label{app:proof_gradient_variance_for_objectives}

\begin{proof}
    Taking $j \sim \text{Unif}(\{1, \dotsc, n\})$, the gradient estimators for the data-dependent terms of $L$ and $L'$ are
\begin{equation}
  \hat g = n \nabla \|\phi(x_j) \theta - \varepsilon_j\|^2_{\Sigma_j^{-1}} = n\phi(x_j)^T \Sigma_j^{-1} (\phi(x_j) \theta - \varepsilon_j)
\end{equation}
and
\begin{equation}
  \hat g' = n \nabla \|\phi(x_j)\theta\|^2_{\Sigma_j^{-1}} = n \phi(x_j)^T \Sigma_j^{-1} \phi(x_j) \theta,
\end{equation}
respectively. Their variances are related as
\begin{align}
  \Var (\hat g)
  &= \Var(n \phi(x_j)^T \Sigma_j^{-1} (\phi(x_j)\theta - \varepsilon_j))  \\
  &= \Var(n\phi(x_j)^T\Sigma_j^{-1}\phi(x_j)\theta) + \Var(n\phi(x_j)^T\Sigma_j^{-1}\varepsilon_j) \nonumber\\ &\quad- 2 \Cov(n\phi(x_j)^T\Sigma_j^{-1}\phi(x_j)\theta,\,\, n \phi(x_j)^T \Sigma_j^{-1}\varepsilon_j)\\
  &= \Var (\hat g') + \Var(n\phi(x_j)^T\Sigma_j^{-1}\varepsilon_j) - 2 \Cov(n\phi(x_j)^T\Sigma_j^{-1}\phi(x_j)\theta,\,\, n\phi(x_j)^T\Sigma_j^{-1} \varepsilon_j)
\end{align}
Evaluating the variance and covariance, we have
\begin{equation}
  \Var \left(n \phi(x_j)^T\Sigma_j^{-1}\varepsilon_j \right) = n\Var(\Phi^T\Sigma^{-1} \varepsilon)
\end{equation}
and
\begin{equation}
 \Cov(n\phi(x_j)^T\Sigma_j^{-1}\phi(x_j)\theta,\,\, n\phi(x_j)^T\Sigma_j^{-1} \varepsilon_j) = n\Cov(\Phi^T\Sigma^{-1}\Phi \theta,\, \Phi^T \Sigma^{-1} \varepsilon),
\end{equation}
and thus
\begin{equation}
  \Var \hat g - \Var \hat g' = n\left[ \Var(\Phi^T\Sigma^{-1}\varepsilon) - 2\Cov(\Phi^T\Sigma^{-1}\Phi \theta, \Phi^T\Sigma^{-1}\varepsilon)\right] \eqqcolon n \Delta.
\end{equation}
\end{proof}

\subsection{Proof of \Cref{eq:effective-dim-condition}}
\label{app:proof_effective_dim_condition}

\begin{proof}
We consider $\Tr\Delta$ for $\theta = \bar{\theta} \sim \mathcal{N}(0, A_{\star}^{-1})$, the optimum of both $L$ and $L'$. From the first order optimality condition, 
\begin{equation}
  \bar{\theta} = H^{-1}(\Phi^T \Sigma^{-1}\varepsilon + A\theta_0).
\end{equation}
Proceeding to rearrange the condition at $\theta=\bar{\theta}$,
\begin{align}
  \Tr{\Delta} 
  &= \Tr\{\E \Phi^T \Sigma^{-1} \varepsilon (\Phi^T \Sigma^{-1} \varepsilon - 2 \Phi^T \Sigma^{-1} \Phi \bar{\theta})^T \} \\
  &= \Tr\{\E \Phi^T \Sigma^{-1} \varepsilon (\Phi^T \Sigma^{-1} \varepsilon - 2 \Phi^T \Sigma^{-1}\Phi H^{-1}(\Phi^T \Sigma^{-1}\varepsilon + A\theta_0))^T \} \\
  &= \Tr\{\E \Phi^T\Sigma^{-1} \varepsilon (\Phi^T\Sigma^{-1} \varepsilon - 2 \Phi^T \Sigma^{-1}\Phi H^{-1}(\Phi^T\Sigma^{-1}\varepsilon + A\E[\theta_0]))^T \}\\
  &= \Tr\{\E \Phi^T \Sigma^{-1}\varepsilon (\Phi^T\Sigma^{-1} \varepsilon - 2 \Phi^T \Sigma^{-1}\Phi H^{-1}\Phi^T \Sigma^{-1}\varepsilon)^T \}, \\
  &= \Tr\{ \Phi^T \Sigma^{-1} \E[\varepsilon \varepsilon^T] (\Sigma^{-1} \Phi - 2 \Sigma^{-1} \Phi H^{-1} \Phi^T \Sigma^{-1} \Phi) \}, \\
  &= \Tr\{ \Phi^T \Sigma^{-1} \Phi(I - 2H^{-1}\Phi^T \Sigma^{-1} \Phi)\}
 \end{align}
 where we substituted in the definition of $\bar{\theta}$, then used that $\varepsilon$ and $\theta_0$ are independent, and that $\E[\theta_0] = 0$, and finally that $\E[\varepsilon\varepsilon^T] = \Sigma$.
 
Writing $M = \Phi^T \Sigma^{-1} \Phi$ and recalling that $H = (M + A)$, we have
\begin{align}
\Tr{\Delta} &= \Tr\{M (I - 2 (M + A)^{-1} M)\}, \\
&= \Tr\{M (I - 2 (A^{-1}M + I)^{-1} A^{-1}M)\}, \\
&= \Tr\{M (I - 2 (I - (A^{-1}M + I)^{-1})\}, \\
&= \Tr\{M (- I + 2(M + A)^{-1} A)\} \\
&= -\Tr\{M\} + 2 \Tr\{M (M + A)^{-1} A\},
\end{align}
where we have used that $(A^{-1}M + I)^{-1} A^{-1}M = I - (A^{-1}M + I)^{-1}$ for the fourth equality. Now consider the isotropic prior case $A = \alpha I$ and recall the effective dimension is written as $\gamma = \Tr\{M (M + A)^{-1}\}$. The above implies $\Tr \Delta > 0$ if and only if $2\alpha \gamma > \Tr \Phi^T \Sigma^{-1} \Phi$.
\end{proof}

\subsection{Analysing Convergence Conditions}
\label{app:large_regulariser_derivation}

To gain some intuition for the condition at convergence, denote by $\lambda_1, \ldots, \lambda_{d^{\prime}}$ the eigenvalues of $M$ (with multiplicity). We can use these to restate the condition as
\begin{equation}
    2 \alpha \gamma=2 \alpha \sum_{j=1}^{d^{\prime}} \frac{\lambda_j}{\lambda_j+\alpha}>\sum_{j=1}^{d^{\prime}} \lambda_j=\operatorname{Tr}\{M\}.
\end{equation}
This formulation of effective dimension gives an interpretation of a soft count of the number of dimensions for which $\lambda_j$ is larger than $\alpha$; in that sense, $\lambda_j$ measures how well determined the corresponding dimension of the weight vector $\theta$ is by the observed data. From here, note that
\begin{equation}
    \frac{2 \alpha \lambda_j}{\lambda_j+\alpha}>\min \left\{\lambda_j, \alpha\right\}.
\end{equation}
and thus it is sufficient for $\operatorname{Tr} \Delta>0$ to hold at convergence that $\alpha>\lambda_j$ for all $j$ (but, of course, not necessary), yielding the intuition that $L^{\prime}$ is preferred when the problem is heavily regularised.

\subsection{Resolving feature scale indeterminacies in the NN Jacobian with the g-prior}\label{app:g-prior-scale-solution}

\textcite{antoran2022linearised} show that for NNs with normalisation layers, the Jacobian features $\phi(\cdot) = \nabla_{\theta}f(\hat{\theta}, \cdot)$ corresponding to each NN layer are scaled arbitrarily. To illustrate this, we divide the NN linearisation point into the concatenation of two weight vectors $\hat{\theta} = [\hat{\theta}_{0}, \hat{\theta}_{1}]$. We assume the layer containing $\hat{\theta}_{0}$ is followed by a normalisation layer, but not that containing $\hat{\theta}_{1}$, which leads to the invariance 
\begin{equation}
    f([k \hat{\theta}_{0}, \hat{\theta}_{1}], \cdot) = f([\hat{\theta}_{0}, \hat{\theta}_{1}], \cdot)
\end{equation}
for all $k>0$.

While $f$ is invariant to this scaling, the Jacobian feature embeddings $\phi(\cdot) = \nabla_\theta f(\hat{\theta}, \cdot)$ are not. We separate the embeddings as 
\begin{equation}
    \phi(x) = [\phi_{0}(\cdot), \phi_{1}(\cdot)] = [\nabla_{\theta_0}f(\hat{\theta}, \cdot), \nabla_{w\theta_1}f(\hat{\theta}, \cdot)].
\end{equation}
\textcite{antoran2022linearised} show that, given a reference pair $([\hat{\theta}_{0}, \hat{\theta}_{1}], [\phi_{0}(x), \phi_{1}(x)])$, and for $\hat{\theta}_{0}$ normalised, scaling $\hat{\theta}_0$ by $k$ results in the pair $([ k \hat{\theta}_{0}, \hat{\theta}_{1}], [k^{-1} \phi_{0}(x), \phi_{1}(x)])$. Thus, using a single prior precision parameter, the regularisation strength applied to the weights multiplying $\phi_{0}(x)$ relative to those multiplying $\phi_{1}(x)$ will increase with $k$. The value of $k$, the scale of the linearisation point, depends on exogenous factors such as learning rate or batch size---and importantly is independent of the data, since it does not affect the output.

One way to resolve this is to assign the weights $\hat{\theta}_0$ and $\hat{\theta}_1$ separate regularisation parameters and learn these using the EM procedure. However, instead, consider using the g-prior normalised features $\phi'$ introduced in \cref{sec:g-prior}, and specifically, the scaling vector corresponding to normalised and non-normalised components $s = [s_{0}, s_{1}]$. For a reference pair $([\hat{\theta}_{0}, \hat{\theta}_{1}], [s_{0}, s_{1}])$ and for $\hat{\theta}_{0}$ normalised, the k-scaled pair is $([ k \hat{\theta}_{0}, \hat{\theta}_{1}]$ and 
\begin{equation*}
   [ \text{diag}(k^{-1} \Phi_{0}^T\mathrm{B} \Phi_{0} k^{-1})^{\odot -\frac{1}{2}},\,\, \text{diag}(\Phi_{1}^T\mathrm{B} \Phi_{1})^{\odot-\frac{1}{2}} ] = [k s_{0}, s_{1}]
\end{equation*}
where $\odot$ denotes an elementwise power. Since the $k$-scaled features are $[k^{-1} \phi_{0}(\cdot), \phi_{1}(\cdot)]$, when applying the g-prior normalisation we recover a feature vector independent of $k$. This resolves the aforementioned pathology.

\section{Experimental Details for \Cref{chapter:bnn_sampling}}

\subsection{MNIST experiments}
\label{app:mnist_experiments}
MNIST $m{=}10$ way classification experiments were performed using the LeNet-style CNN architectures of increasing size employed by \textcite{antoran2022linearised}: ``LeNetSmall'' ($d'{=}14634$), ``LeNet'' ($d'{=}29226$) and ``LeNetBig'' ($d'{=}46024$). We note that these models contain batch normalisation layers. Each model was trained with using SGD with momentum of 0.9 for 90 epochs with a learning rate drop of a factor of 10 every 30 epochs. The MNIST dataset was downloaded from \texttt{PyTorch torchvision}. We employ standard per-channel mean and std-dev standardisation preprocessing and two pixel shift and crop data augmentation. For posterior mode optimisation and sampling, we do not perform data augmentation as to avoid cold posterior effects \citep{izmailov2021what}.
The details of our SGD approaches to convex optimisation for obtaining posterior modes and samples are as follows
\begin{itemize}
    \item \textbf{Posterior mode optimisation}: The linearised NN weights are trained using SGD with a Nesterov momentum coefficient of 0.9, and batch size 1000 for 40 epochs. We clip gradients to a maximum norm of $1$. We use an initial learning rate of $1e-2$ when using standard isotropic or layerwise Gaussian priors, and $1$ for the g-prior. We employ a linear decay schedule that reduces the lr by a factor of 330 over the first 75$\%$ of the training procedure and holds it constant afterwards.
    \item \textbf{Sampling}: We optimise $32$ samples in parallel using SGD with Nesterov momentum $({=}0.9)$ and a batch size of 1000 for 20 epochs. For standard Gaussian priors (isotropic and layerwise), we use a learning rate of $2e{-}1$, whereas for the g-prior, we find a higher learning rate of $200$ to work best. When using the KFAC baseline, we draw posterior samples through eigendecomposition, following \textcite{daxberger2021laplace}.
    % We keep an exponential moving Polyak average, with an update size of $0.01$ every train step. 
\end{itemize}
\textbf{Hyperparameter optimisation}: We tuned the learning rate, decay schedule and gradient clipping strength using a rough grid search over multiple orders of magnitude. We chose the settings that reached the lowest loss values. These can be evaluated with just the train set. We chose the largest batch size that could accommodate optimising $32$ samples in parallel on a single hardware accelerator. We note that posterior mode and sample optimisation converge in less than half of the total epochs we use for their optimisation. The numbers of epochs chosen were set to be large enough to ensure convergence and not tuned. A decrease in computational cost can likely be achieved by stopping sample optimisation earlier. 

\paragraph{Baseline methods.} \ For the comparison of learning a single precision hyperparameter and layerwise hyperparameters in \Cref{fig:mnist_preds_and_EM}, we extend the M-step update to as $\alpha_{l} = \nicefrac{\gamma_{l}}{\|\hat{\theta}_{\text{lin}, l}\|_{2}^2}$ where $l$ indexes each layer's attributes, as done in \textcite{Mackay1992Thesis,Tipping2001sparse}. For the MAP, diagonal covariance and KFAC covariance baselines,
% as well as the additional baselines employed in \cref{app:add_experiments},
we use the same pre-trained models when possible (i.e. not for the ensembles or dropout baselines). Since all baselines share the same linearisation point, they also share the same mean predictions. Differences in performance among baselines are thus only due to differences in uncertainty estimation. The diagonal approximation to the covariance is constructed by first computing the diagonal of the Hessian $M$ and the inverting it. For the KFAC covariance approximation, we exploit the equivalency between the Generalised Gauss Newton matrix (i.e. the Hessian of the linear model $h$) and the Fisher information matrix for exponential family likelihoods (i.e. the categorical). This allows us to formulate the Hessian as an expectation of likelihood gradients, which in turn we approximate using a single sample per training observation, as in \textcite{daxberger2021laplace}. 

For completeness, we also state the probit approximation for sampled predictive posteriors over logits. For input $x$ and samples $\zeta_{i}, \dotsc, \zeta_{k}$, the predictive probability for class $i \in \|1, \dotsc, m\|$ is given by 
\begin{align*}
    \text{softmax}\left(f(\hat{\theta}_{\text{nn}}, x) \odot (1 + \frac{\pi}{2k} \sum_{j < k} (\phi(x)\zeta_{j})^{\odot 2} )^{\odot -0.5} \right)_{i}. \label{eq:probit_approx}
\end{align*}

\subsection{CIFAR100 classification}
\label{app:cifar100_experiments}
CIFAR100 $m{=}100$ way classification experiments were performed using ResNet18 models ($d'\approx 11M$) with specific architecture details matching the \href{https://pytorch.org/vision/main/models.html}{\texttt{PyTorch torchvision implementation}}. We train these models using SGD with momentum of 0.9 for 300 epochs. The starting lr is 0.1 and we reduce it by a factor of 10 every 100 epochs. The CIFAR100 dataset was also downloaded using \texttt{torchvision} and our data preprocessing and augmentation also follow the default implementation from this library. For posterior mode optimisation and sampling, we do not perform data augmentation since doing so would result in non-iid  observations and an intractable likelihood.
The SGD details used to solve the convex optimisation problems required for obtaining posterior modes and drawing samples are as follows
\begin{itemize}
    \item \textbf{Posterior mode optimisation}: The linearised NN weights are trained using SGD with Nesterov momentum $({=}0.9)$ and a batch size of 2000 for 40 epochs. We employ a linear decay learning rate schedule with an initial learning rate of $1e{-}1$. It is decreased by a factor of $330$ over the first $75\%$ of training, and then held constant. We also employ gradient clipping with maximum norm$=0.1$.
    \item \textbf{Sampling}: We optimise $6$ samples in parallel using SGD with Nesterov momentum $({=}0.9)$ and a batch size 100 for 20 epochs. All other details match those of posterior mode optimisation. 
\end{itemize}
Upon convergence of the EM algorithm, we draw $64$ further samples using the optimal prior precision by following the optimisation procedure described above. We initialise these samples at prior samples drawn with the optimised prior precision. When using the KFAC baseline, we draw posterior samples through eigendecomposition \citep{daxberger2021laplace}.

\textbf{Hyperparameter optimisation} We tuned the learning rate, decay rate and gradient clipping strength using a rough grid search over orders of magnitude. This search only requires access to training data, since a lower loss value is always preferable when optimising samples. We chose the largest batch size for which we could simultaneously optimise $6$ samples in parallel on a single hardware accelerator. Similarly to the MNIST experiments, we did not optimise the number of optimisation epochs and instead chose large values that would ensure convergence. It is likely that our EM iteration can be sped up by decreasing the duration of the convex optimisation routines.

\paragraph{EM iteration with KFAC} When using the KFAC approximation, we learn layer-wise regularisation strength parameters with the EM algorithm. We use the same procedure as for our sampling approximation, given in \Cref{alg:sampling_classification}, but we do not use Mackay's fixed point iteration for the M step. When using the KFAC approximation to the posterior covariance, we find this iteration to suffer from strong oscillations of the regularisation parameter. These lead to slow convergence and even divergence in some settings.
We instead update the regularisation parameters in the M step by applying gradient based optimisation to the evidence lower bound given in \Cref{def:fixed_cov_elbo}. The KFAC approximation allows for fast log-determinant evaluation, making repeated evaluations of this objective feasable. 

Details for baselines and hyperparameters not mentioned explicitly in this subsection match those given for MNIST in the previous subsection. 

\subsubsection{Efficient $\kappa$-adic sampling}
\textcite{osband2022evaluating,Osband2021epistemic} introduced \textit{dyadic} test input sampling ($\kappa = 2)$ as a practical way of evaluating joint predictions in discriminative tasks. This method samples $\kappa=2$ random anchor points from the test dataset, and then randomly resamples them to create a batch of size $\tau=10$. Test log-likelihood is evaluated jointly for each batch as
\begin{gather*}
   \log  \int \exp\left(\sum_{i \leq \tau}\ell(y_{i}, f(\theta, x_{i})) \right) \, d\Pi,
\end{gather*}
for $f$ the model being evaluated and $\Pi$ its posterior distribution over model parameters. This quantity can be estimated with posterior samples $\zeta_{1}, \dotsc, \zeta_{k} \sim \Pi$ as
\begin{gather*}
   \log  \frac{1}{k} \sum_{j \leq k} \exp\left(\sum_{i \leq \tau}\ell(y_{i}, f(\zeta_{j}, x_{i})) \right).
\end{gather*}
We extend this evaluation approach to larger $\kappa$ and $\tau$ values without increasing computational cost. 
We randomly sample $\kappa$ integers $\{b_1, \ldots, b_{\kappa}\}$ such that they sum to $\tau$, i.e $\sum_i^{\kappa} k_i = \tau$. The joint log-likelihood over the batch of size $\tau$ with $\kappa$ unique datapoints can then be estimated as 
\begin{gather*}
    \log  \frac{1}{k} \sum_{j \leq k} \exp\left(\sum_{l \leq \kappa} b_{l} \ell(y_{l}, f(\zeta_{j}, x_{l})) \right).
\end{gather*}
where the inner sum is over the $\kappa$ distinct elements in the batch instead of the ``total batch size'' $\tau$.
This is equivalent to the formulation proposed in \textcite{osband2022evaluating} for dyadic sampling, when $\kappa=2$ and $\tau = 10$. We note that it is not possible to achieve \textit{augmented dyadic} sampling, as described in \textcite{Osband2021epistemic}, with this approach. However the authors mention that there is not a significant difference in the relative performance of methods when using augmented dyadic sampling compared to regular dyadic sampling. We introduce a final step however, which is to repeat the computation for multiple shuffles (10) of the test dataset. This eliminates variance in our results from the choice of the $\kappa$ observations which get grouped together in each batch. 

\subsection{Imagenet classification}
\label{app:imagenet_experiments}
ImageNet $m{=}1000$ way classification experiments were performed using ResNet50 models ($d'\approx 25M$) with specific architecture details matching the \href{https://pytorch.org/vision/main/models.html}{\texttt{PyTorch torchvision implementation}}. We train these models using the default \texttt{torchvision} settings. We use SGD with momentum of 0.9 for 90 epochs. The starting lr is 0.1 and we reduce it by a factor of 10 every 30 epochs. The ImageNet dataset was downloaded from the official website \href{https://www.image-net.org/}{\texttt{here}} and our data preprocessing and augmentation follow the default implementation from the \texttt{torchvision} library. For posterior mode optimisation and sampling, we do not perform data augmentation.
The SGD details used to solve the convex optimisation problems required for obtaining posterior modes and drawing samples are as follows
\begin{itemize}
    \item \textbf{Posterior mode optimisation}: The linearised NN weights are trained using SGD with Nesterov momentum $({=}0.9)$ and a batch size of 240 for 12 epochs. We employ a linear decay learning rate schedule with an initial learning rate of $5$. It is decreased by a factor of $33$ over the first $95\%$ of training, and then held constant. We also employ gradient clipping with maximum norm$=1$. This procedure takes around 12 hours on a TPU-v3 accelerator.
    \item \textbf{Sampling}: We optimise $6$ samples in parallel using SGD with Nesterov momentum $({=}0.9)$ and a batch size 24 for 1 epoch. We employ a linear decay learning rate schedule with an initial learning rate of $100$. It is decreased by a factor of $3330$ over the first $95\%$ of training, and then held constant. This procedure takes around 8 hours on a TPU-v3 accelerator. All other details match those of posterior mode optimisation.
\end{itemize}
Upon convergence of the EM algorithm, we draw $90$ further samples using the optimal prior precision by following the optimisation procedure for sampling described above. We initialise these samples at prior samples drawn with the optimised prior precision.

\textbf{Hyperparameter optimisation}: We tuned the learning rate, decay rate and gradient clipping strength using a rough grid search over orders of magnitude. We also chose the largest batch size that for which we could simultaneously optimise $6$ samples in parallel on a single hardware accelerator.

Details for baselines and hyperparameters not mentioned explicitly in this subsection match those given for CIFAR100 and MNIST in the previous subsections. 

\subsection{Tomographic reconstruction}\label{app:DIP_experimental_details}

\paragraph{Problem setup} Tomographic reconstruction consists in solving a linear inverse problem in imaging where we observe a set of measurements $y \in \R^{m}$, which we assume to be generated as $y \,{=}\, U x^* + \eta$ for $U \in \R^{m \times d}$ the discrete Radon transform, $x^* \in \R^{d}$ the image to reconstruct and $\eta \sim \mathcal{N}(0, I)$ random noise. We have $m \ll d$, making the problem underconstrained. 
Our specific tomographic reconstruction task closely follows the one from \textcite{barbano2021deep}. We perform a sparse-view reconstruction of an image of a slice of a walnut from a sub-sampled set of measurements. Specifically, from the full measurement set of \textcite{der2019cone}, which containing scans at 1200 equidistant angles over $[0, 360^\circ)$, we choose our measurement set by subsampling angles by a factor of either 10x or 20x, leading to measurements of size $m=15360$ or $m=7680$.
As in \textcite{barbano2021deep, antoran2022tomography,barbano2023painless}, we reduce the original 3D scan geometry to the 2D slice of interest by selecting the relevant subset of measurement pixels. We assemble the Radon operator $U$ as a sparse matrix taking in an image of resolution $(501$px$)^2$ and outputting a measurement tensor coherent with the described geometry. 

\paragraph{Methods}

To provide a reconstruction, we use the Deep Image prior \citep{Ulyanov2020deep} which trains the parameters $w \in \R^{d'}$ of a fully convolutional U-Net autoencoder $f : \R^{d'} \to \R^{d}$, where the input is fixed and thus absent from our notation, until a satisfactory reconstruction $f(\hat{\theta}_{\text{nn}})$ is obtained.  The U-Net network architecture is the one proposed in \textcite{baguer2020computed}.
The optimisation of the U-Net parameters follows \textcite{barbano2021deep}, although we note that faster optimisation strategies exist \citep{barbano2023painless}.
To estimate the uncertainty in this reconstruction, we linearise the U-Net around $\hat{\theta}_{\text{nn}}$, as described in \Cref{def:lla}. 
This leaves us with a model affine in the parameters and with design matrix $U\Phi \in \R^{m \times d'}$. We may now proceed with linear model inference.
While \textcite{antoran2022tomography} use the traditional EM iteration described in \Cref{def:first_order_em}, we use the sample-based one from \Cref{sec:sample_em}. 
For evaluation, we use the non-sparse reconstruction (using 1200 angles) provided by \textcite{der2019cone} as the ground truth image $x^*$.
To evaluate joint log-likelihoods we estimate the predictive covariance matrix for patches of neighbouring pixels using samples. Covariance matrix estimates from samples are known to be unreliable. We use the stabilised formulation of \textcite{Maddox2019Simple}: $\hat\Sigma = \frac{1}{2k} \left[\sum_{j=1}^k \text{diag}(\hat{x}_{j}^{2}) + \hat{x}_{j} \hat{x}_{j}^{T}\right]$ for $(\hat{x}_{j})_{j=1}^k$ samples from the predictive posterior over a patch.
We then construct predictive distributions over pixels as $\mathcal{N}(f(\hat{\theta}_{\text{nn}}), \hat\Sigma)$.

\paragraph{Hyperparameters} We employ a low CG tolerance of $1e-3$ and a maximum number of iterations of $150$, which is never reached in practise as the error tolerance level is always hit in less steps. Our dropout baseline follows \textcite{antoran2022tomography} in placing a dropout layer after every 2d convolution. The dropout rate is set to 0.05. 